\section*{अध्याय \ref{ch:g6:pure-substances-and-mixtures} --- शुद्ध पदार्थ और मिश्रण}

\begin{solution}{exo:g6:pure-substances-and-mixtures:1}
\omterm{def:g6:pure-substances-and-mixtures:pure-substance}{शुद्ध पदार्थ}: आसुत जल, शुद्ध लोहे का टुकड़ा, चीनी का
क्रिस्टल। \omterm{def:g6:pure-substances-and-mixtures:mixture}{मिश्रण}: \omterm{def:g4:air-a-mixture-of-gases:air}{वायु}, समुद्र का पानी, दूध।
\end{solution}

\begin{solution}{exo:g6:pure-substances-and-mixtures:2}
समांगी: नमकीन पानी, नल का पानी। विषमांगी: ग्रेनाइट, पानी पर
तेल, मूसली।
\end{solution}

\begin{solution}{exo:g6:pure-substances-and-mixtures:3}
किसी पदार्थ का एक विशेष प्रकार, जिसका अपना नाम और अपने गुण होते हैं: जैसे
पानी, नमक, एथेनॉल (और भी: चीनी, ऑक्सीजन, लोहा)।
\end{solution}

\begin{solution}{exo:g6:pure-substances-and-mixtures:4}
निश्चित क्वथन ताप से \omterm{def:g6:pure-substances-and-mixtures:pure-substance}{शुद्ध पदार्थ} का संकेत मिलता है; \qty{78}{\celsius}
एथेनॉल का क्वथन ताप है।
\end{solution}

\begin{solution}{exo:g6:pure-substances-and-mixtures:5}
उसमें अनेक \omterm{def:g6:pure-substances-and-mixtures:species}{रासायनिक स्पीशीज़} हैं (पानी, शर्कराएँ, अम्ल, विटामिन,
सुगंध): वह एक \omterm{def:g6:pure-substances-and-mixtures:mixture}{मिश्रण} है। डिब्बे पर लिखे ``शुद्ध'' का अर्थ केवल इतना है कि
कुछ मिलाया नहीं गया।
\end{solution}

\begin{solution}{exo:g6:pure-substances-and-mixtures:6}
नहीं: \omterm{def:g6:pure-substances-and-mixtures:pure-substance}{शुद्ध पदार्थ} एक निश्चित ताप पर उबलता है। यह एक \omterm{def:g6:pure-substances-and-mixtures:mixture}{मिश्रण} है,
शायद ऐसा पानी जिसमें कुछ घुला है, जैसे नमकीन पानी।
\end{solution}

\begin{solution}{exo:g6:pure-substances-and-mixtures:7}
छन्ना गंदले पानी को (कण पेपर पर रह जाते हैं) और गूदे वाले रस को
(गूदा रह जाता है) अलग कर सकता है। वह नमकीन पानी को (नमक घुला हुआ है) अलग
नहीं कर सकता, न पानी से तेल को (दोनों द्रव निकल जाते हैं; उन्हें ऊपर की परत
उँडेलकर अलग किया जाता है)।
\end{solution}

\begin{solution}{exo:g6:pure-substances-and-mixtures:8}
\qty{1}{cm^3} के लिए $78.7 \div 10 = \qty{7.87}{g}$: लोहा।
\end{solution}

\begin{solution}{exo:g6:pure-substances-and-mixtures:9}
$\frac{357}{478} \approx 0.747$: लगभग \qty{75}{\percent}।
\end{solution}

\begin{solution}{exo:g6:pure-substances-and-mixtures:10}
$2 \times 35 = \qty{70}{g}$ लवण। प्रतिशत: $\frac{35}{1000} =
\qty{3.5}{\percent}$।
\end{solution}

\begin{solution}{exo:g6:pure-substances-and-mixtures:11}
उनके क्वथन ताप मापो (उबलते समय हर एक स्थिर रहना चाहिए, और दोनों
बराबर होने चाहिए), और हर एक का समान आयतन तौलो (बराबर
द्रव्यमान)। दो बराबर स्थिरांक, दोनों स्थिर, एक ही \omterm{def:g6:pure-substances-and-mixtures:pure-substance}{शुद्ध पदार्थ} की ओर
इशारा करते हैं।
\end{solution}

\begin{solution}{exo:g6:pure-substances-and-mixtures:12}
$\frac{10}{90 + 10} = \qty{10}{\percent}$ एथेनॉल। यह एक \omterm{def:g6:pure-substances-and-mixtures:mixture}{मिश्रण} है: इसका
कोई निश्चित क्वथन ताप नहीं होता और यह सारणी में कहीं नहीं है, जिसमें केवल
\omterm{def:g6:pure-substances-and-mixtures:pure-substance}{शुद्ध पदार्थ} दिए गए हैं।
\end{solution}

\begin{solution}{pb:g6:pure-substances-and-mixtures:1}
\textbf{1.} नहीं। \omterm{def:g6:pure-substances-and-mixtures:pure-substance}{शुद्ध पदार्थ} और \omterm{def:g6:pure-substances-and-mixtures:homogeneous}{समांगी मिश्रण} बिल्कुल एक जैसे दिख सकते
हैं: कुछ मापना पड़ता है।

\textbf{2.} समांगी: वह स्वच्छ है और उसके हिस्से दिखाई नहीं देते।

\textbf{3.} कोई स्थिरांक: ऐसा क्वथन ताप जो द्रव के उबलते समय स्थिर
रहे (या किसी ज्ञात आयतन का द्रव्यमान, सारणी से तुलना करके)।

\textbf{4.} A और B, जिनके क्वथन ताप स्थिर रहते हैं, \omterm{def:g6:pure-substances-and-mixtures:pure-substance}{शुद्ध
पदार्थ} हैं। C एक \omterm{def:g6:pure-substances-and-mixtures:mixture}{मिश्रण} है।

\textbf{5.} A \qty{100}{\celsius} पर उबलता है: पानी। B
\qty{78}{\celsius} पर उबलता है: एथेनॉल।

\textbf{6.} A: $49.8 \div 50.0 = \qty{0.996}{g}$। B: $39.5 \div 50.0 =
\qty{0.790}{g}$। C: $51.2 \div 50.0 = \qty{1.024}{g}$।

\textbf{7.} हाँ: पानी के लिए लगभग \qty{1.0}{g}, एथेनॉल के लिए
\qty{0.79}{g}।

\textbf{8.} जैसे-जैसे पानी उबलकर निकलता है, नमक पीछे रह जाता है, इसलिए बचे
हुए द्रव में नमक बढ़ता जाता है; तब वह और ऊँचे ताप पर
उबलता है।

\textbf{9.} द्रव C में एक घुला हुआ ठोस है: वह एक विलयन है, यानी एक
\omterm{def:g6:pure-substances-and-mixtures:mixture}{मिश्रण}, जैसा उसके उबलने से पहले ही पता चल गया था।

\textbf{10.} $100 \times 1.024 = \qty{102.4}{g}$; $\frac{3.5}{102.4}
\approx 0.034$: द्रव्यमान का लगभग \qty{3.4}{\percent} नमक है।

\textbf{11.} $\frac{35}{1000} = \qty{3.5}{\percent}$: प्रश्न 10 के
बहुत क़रीब।

\textbf{12.} \qty{100}{mL} में \qty{3.5}{g}, इसलिए $3.5 \times 10 =
\qty{35}{g}$ नमक प्रति लीटर। अपने खारेपन और उबलने के ढंग के साथ C समुद्र
का पानी है; A आसुत जल है और B एथेनॉल।
\end{solution}
